Considere la ecuación de Laplace $u_{xx}+u_{yy}=0$ en el rectángulo
$0<x<a$, $0<y<b$, con $a,b>0$ y lados verticales aislados:
\[
u_x(0,y)=u_x(a,y)=0,\qquad 0<y<b.
\]
Resuelva los tres apartados en el orden indicado.

Los datos $f,g$ pertenecen a $C^2[0,a]$. Se busca una función continua
en el rectángulo cerrado, armónica en su interior y con derivadas hasta
orden dos continuas junto a los lados verticales para $0<y<b$.
Las condiciones Neumann se imponen en esos lados abiertos. No se impone
regularidad de las derivadas en los cuatro vértices; exigirla añadiría
compatibilidades a los datos.

\begin{enumerate}[label=(\alph*),ref=\alph*,leftmargin=*]
\item Obtenga por separación de variables la solución con $u(x,0)=f(x)$ y
$u(x,b)=0$. Incluya todos los modos espaciales admisibles.
\item Utilice el resultado del apartado anterior para resolver el problema con
$u(x,0)=f(x)$ y $u(x,b)=g(x)$.
\item Aplique el apartado anterior a $a=\pi$, $b=1$,
$f(x)=\cos x-\cos3x$ y $g(x)=\cos2x$.
\end{enumerate}
